%% kapa: exam=1 round=35 year=2024 period=2 form=A subject=appraisal_law subject_ko=감정평가관계법규 seq=36 number=36
%% source: Q-Net 감정평가사 기출문제 (cst00309) artlSeq=5213229 file=raw/1008/5213229/2교시 문제지 원본.pdf
%% conversion: kordoc 4.24.51 -> tools/gen1.py (auto); review=proofread-sonnet; flags=none
\begin{question}{36}
동산ㆍ채권 등의 담보에 관한 법령상 담보등기에 관한 설명으로 옳은 것은?
\begin{choices}
  \item 장래에 취득할 동산은 특정할 수 있는 경우에도 이를 목적으로 담보등기를 할 수 없다.
  \item 등기명의인 표시의 변경의 등기는 등기명의인 단독으로 신청할 수 있다.
  \item 담보권자가 담보권의 존속기간을 갱신하려면 그 존속기간 만료 전후 1개월 내에 연장 등기를 신청하여야 한다.
  \item 포괄승계로 인한 등기는 등기권리자 또는 등기의무자 단독으로 신청할 수 있다.
  \item 담보목적물인 동산이 멸실된 경우 그 말소등기의 신청은 담보권설정자가 하여야 한다.
\end{choices}
\end{question}
\answer{2}
